What drives the improvements in downstream performance of multi-token prediction models on all of the tasks we have considered? By conducting toy experiments on controlled training datasets and evaluation tasks, we demonstrate that multi-token prediction leads to \emph{qualitative changes in model capabilities and generalization behaviors}.
In particular, Section~\ref{sect:induction} shows that for small model sizes, \emph{induction capability}---as discussed by \citet{olsson2022context}---either only forms when using multi-token prediction as training loss, or it is vastly improved by it. Moreover, Section~\ref{sect:algorithmic} shows that multi-token prediction improves generalization on an arithmetic task, even more so than tripling model size.

\subsection{Induction capability}
\label{sect:induction}
\begin{figure}[t!]
    \centering
    \includegraphics[width=0.8\linewidth]{img/induction.pdf}
    \caption{\textbf{Induction capability of $n$-token prediction models.} 
     Shown is accuracy on the second token of two token names that have already been mentioned previously. Shown are numbers for models trained with a next-token and a 2-token prediction loss, respectively, with two independent runs each. The lines denote per-loss averages. For small model sizes, next-token prediction models learn practically no or significantly worse induction capability than 2-token prediction models, with their disadvantage disappearing at the size of 100M nonembedding parameters.
}

    \label{fig:induction_size}
\end{figure}

Induction describes a simple pattern of reasoning that completes partial patterns by their most recent continuation \citep{olsson2022context}. In other words, if a sentence contains ``AB'' and later mentions ``A'', induction is the prediction that the continuation is ``B''. We design a setup to measure induction capability in a controlled way. Training small models of sizes 1M to 1B nonembedding parameters on a dataset of children stories, we measure induction capability by means of an adapted test set: in 100 stories from the original test split, we replace the character names by randomly generated names that consist of two tokens with the tokenizer we employ.
Predicting the first of these two tokens is linked to the semantics of the preceding text, while predicting the second token of each name's occurrence after it has been mentioned at least once can be seen as a pure induction task. In our experiments, we train for up to 90 epochs and perform early stopping with respect to the test metric (i.e. we allow an epoch oracle). Figure~\ref{fig:induction_size} reports induction capability as measured by accuracy on the names' second tokens in relation to model size for two runs with different seeds.

We find that 2-token prediction loss leads to a vastly improved formation of induction capability for models of size 30M nonembedding parameters and below, with their advantage disappearing for sizes of 100M nonembedding parameters and above.\footnote{Note that a perfect score is not reachable in this benchmark as some of the tokens in the names in the evaluation dataset never appear in the training data, and in our architecture, embedding and unembedding parameters are not linked.}
We interpret this finding as follows: multi-token prediction losses help models to learn transferring information across sequence positions, which lends itself to the formation of induction heads and other in-context learning mechanisms. However, once induction capability has been formed, these \emph{learned features} transform induction into a task that can be solved \emph{locally} at the current token and learned with next-token prediction alone. From this point on, multi-token prediction actually hurts on this restricted benchmark---but we surmise that there are higher forms of in-context reasoning to which it further contributes, as evidenced by the results in Section~\ref{sect:model-scaling}. In Figure~\ref{fig:induction_b3g}, we provide evidence for this explanation: replacing the children stories dataset by a higher-quality 9:1 mix of a books dataset with the children stories, we enforce the formation of induction capability early in training by means of the dataset alone. By consequence, except for the two smallest model sizes, the advantage of multi-token prediction on the task disappears: feature learning  of induction features has converted the task into a pure next-token prediction task.

\subsection{Algorithmic reasoning}
\label{sect:algorithmic}
\begin{figure}[t!]
    \centering
    \includegraphics[width=0.95\linewidth]{img/poly_100M_pause_0.pdf}
    \caption{\textbf{Accuracy on a polynomial arithmetic task with varying number of operations per expression.} Training with multi-token prediction losses increases accuracy across task difficulties. In particular, it also significantly improves out-of-domain generalization performance, albeit at a low absolute level. Tripling the model size, on the other hand, has a considerably smaller effect than replacing next-token prediction with multi-token prediction loss (Figure~\ref{fig:poly_scale}). Shown are two independent runs per configuration with 100M parameter models.
}
    \label{fig:poly_0}
\end{figure}

Algorithmic reasoning tasks allow to measure more involved forms of in-context reasoning than induction alone. We train and evaluate models on a task on polynomial arithmetic in the ring $\mathbb{F}_7[X]/(X^5)$ with unary negation, addition, multiplication and composition of polynomials as operations. The coefficients of the operands and the operators are sampled uniformly. %
The task is to return the coefficients of the polynomials corresponding to the resulting expressions. The number $m$ of operations contained in the expressions is selected uniformly from the range from 1 to 5 at training time, and can be used to adjust the difficulty of both in-domain ($m \leq 5$) and out-of-domain ($m > 5$) generalization evaluations. The evaluations are conducted with greedy sampling on a fixed test set of 2000 samples per number of operations. We train models of two small sizes with 30M and 100M nonembedding parameters, respectively. This simulates the conditions of large language models trained on massive text corpora which are likewise under-parameterized and unable to memorize their entire training datasets.

Multi-token prediction improves algorithmic reasoning capabilities as measured by this task across task difficulties (Figure~\ref{fig:poly_0}). In particular, it leads to impressive gains in out-of-distribution generalization, despite the low absolute numbers. Increasing the model size from 30M to 100M parameters, on the other hand, does not improve evaluation accuracy as much as replacing next-token prediction by multi-token prediction does (Figure~\ref{fig:poly_scale}). In Appendix~\ref{app:poly}, we furthermore show that multi-token prediction models retain their advantage over next-token prediction models on this task when trained and evaluated with \emph{pause tokens} \citep{goyal2023think}.
